\lecture{9.8}{Missing Data: Traditional Statistical View}
\label{lec:9-8}

\subsection{Target Estimands With Incomplete Data}

Missing data problems begin with a target estimand defined for a complete-data
distribution. Examples include
\[
  \psi_1=\E(Y),
\]
\[
  \psi_2=\E(Y\given X=x),
\]
and, under causal assumptions,
\[
  \psi_3=\E\{\E(Y\given X,T=1)-\E(Y\given X,T=0)\}.
\]
These estimands are functions of variables such as $(X,T,Y)$ that may not be
fully observed for every sampled unit.

\subsection{Missingness Indicators}

Let $V$ be a variable that may be missing. Define
\[
  R_V=
  \begin{cases}
  1, & V \text{ is observed},\\
  0, & V \text{ is missing}.
  \end{cases}
\]
The observed version is
\[
  V^*=
  \begin{cases}
  V, & R_V=1,\\
  ?, & R_V=0.
  \end{cases}
\]
If $X,T,Y$ can all be missing, then the observed data include
\[
  X^*,T^*,Y^*,R_X,R_T,R_Y.
\]

\begin{figure}[H]
\centering
\begin{tikzpicture}[
  box/.style={draw, rounded corners=2pt, minimum width=1.1cm, minimum height=.7cm},
  arr/.style={-{Latex[length=2.2mm]}, thick},
  node distance=1.0cm
]
\node[box] (V) {$V$};
\node[box, below=of V] (R) {$R_V$};
\node[box, right=of R] (Vstar) {$V^*$};
\draw[arr] (V) -- (Vstar);
\draw[arr] (R) -- (Vstar);
\end{tikzpicture}
\caption{The observed variable $V^*$ is determined by the full value $V$ and the missingness indicator $R_V$.}
\label{fig:ch9-missing-indicator}
\end{figure}

\subsection{Complete-Case Analysis}

Complete-case analysis discards any row with missing values in variables needed
for the analysis. It is simple, but it changes the target population from all
sampled units to units with complete records. It is unbiased for some estimands
only under strong missingness conditions, such as missingness independent of the
variables relevant to the estimand.

\subsection{Imputation}

Imputation replaces missing values with predicted or sampled values. Single
imputation fills in one value; multiple imputation creates several completed
data sets and combines results. Imputation can improve efficiency relative to
complete-case analysis, but it depends on assumptions about the missingness
mechanism and the imputation model.

Imputation should not be viewed as a purely technical repair. It encodes
assumptions about what the missing values would have been.

\subsection{Efficiency Loss}

Ignoring incomplete rows reduces sample size and can discard partially observed
information. For example, a row with $X$ and $T$ observed but $Y$ missing still
contains information about the distribution of covariates and treatment. A row
with $Y$ observed but one covariate missing still contains outcome information.
Efficient methods use all observed components compatible with the identifying
assumptions.

\subsection{Bias From Missingness}

Bias occurs when complete cases differ systematically from incomplete cases in
ways related to the target estimand. For example, if people with severe disease
are less likely to report an outcome, the observed outcome mean among
respondents can underestimate the true mean. If missingness depends on the
unobserved value itself, naive analyses can be badly biased.

\subsection{Sources of Missing Data}

Common sources include:
\begin{itemize}
  \item random data-entry failures;
  \item survey nonresponse to sensitive questions;
  \item variables introduced partway through a study;
  \item measurement limits, such as values too fast or too small to record;
  \item dropout or loss to follow-up;
  \item clinician decisions about whether to order a test.
\end{itemize}
Different sources imply different assumptions. A randomly skipped data-entry
field is not the same as a refused sensitive survey question or a diagnostic
test ordered only for high-risk patients.

\begin{keyidea}
Missing data analysis is an identification problem. Before choosing complete
cases, imputation, or weighting, specify the target complete-data estimand and
the assumptions linking observed records to missing values.
\end{keyidea}
